\documentclass[11pt,a4paper]{article}
\usepackage[utf8]{inputenc}
\usepackage[T1]{fontenc}
\usepackage[english]{babel}
\usepackage{amsmath,amsthm,amssymb}
\usepackage{mathtools}
\usepackage[margin=2.3cm]{geometry}
\usepackage{hyperref}
\usepackage{enumitem}
\usepackage{xcolor}
\usepackage{tikz}
\usepackage{booktabs}

% --- Theorem environments ---
\theoremstyle{plain}
\newtheorem{theorem}{Theorem}[section]
\newtheorem{proposition}[theorem]{Proposition}
\newtheorem{lemma}[theorem]{Lemma}
\newtheorem{corollary}[theorem]{Corollary}
\newtheorem{conjecture}[theorem]{Conjecture}
\newtheorem{blueprint}[theorem]{Conditional Blueprint}
\newtheorem{milestone}[theorem]{Milestone}
\newtheorem{hypothesis}[theorem]{Hypothesis}

\theoremstyle{definition}
\newtheorem{definition}[theorem]{Definition}
\newtheorem{remark}[theorem]{Remark}
\newtheorem{example}[theorem]{Example}

% --- Operators ---
\DeclareMathOperator{\spec}{spec}
\DeclareMathOperator{\Spec}{Spec}
\DeclareMathOperator{\Dom}{Dom}
\DeclareMathOperator{\Gal}{Gal}
\DeclareMathOperator{\Tr}{Tr}
\DeclareMathOperator{\Res}{Res}
\DeclareMathOperator{\vol}{vol}
\DeclareMathOperator{\id}{id}
\DeclareMathOperator{\sgn}{sgn}
\DeclareMathOperator{\Hom}{Hom}
\DeclareMathOperator{\End}{End}
\DeclareMathOperator{\Aut}{Aut}
\DeclareMathOperator{\Rep}{Rep}

% --- Shortcuts ---
\newcommand{\C}{\mathbb{C}}
\newcommand{\R}{\mathbb{R}}
\newcommand{\Q}{\mathbb{Q}}
\newcommand{\Z}{\mathbb{Z}}
\newcommand{\N}{\mathbb{N}}
\newcommand{\A}{\mathbb{A}}
\newcommand{\F}{\mathbb{F}}
\newcommand{\HH}{\mathcal{H}}
\newcommand{\V}{\mathcal{V}}
\renewcommand{\S}{\mathcal{S}}
\newcommand{\E}{\mathcal{E}}
\newcommand{\D}{\mathcal{D}}
\newcommand{\K}{\mathcal{K}}
\newcommand{\M}{\mathcal{M}}
\newcommand{\T}{\mathcal{T}}
\newcommand{\chifn}{\chi}
\newcommand{\gap}{\operatorname{gap}}
\renewcommand{\Re}{\operatorname{Re}}
\renewcommand{\Im}{\operatorname{Im}}

% --- Metadata ---
\title{Functional Stability Theory as a Comparative Framework for\\
the Spectrum-Zero Identification:\\
Classification, Three-Route Map, and Reduction Blueprints\\
\vspace{0.3em}
{\large [Extended Draft v0.8 --- Three-Lemma Endgame for Route~B]}}
\author{Lukas Geiger\\
\small Independent Researcher, Bernau}
\date{April 20, 2026}

\begin{document}

\maketitle

\begin{abstract}
This paper presents the Functional Stability Theory (FST) axiom
package for the spectrum-zero identification underlying the Riemann
Hypothesis (RH) and its Dirichlet generalisation, together with
a \emph{conditional endgame} for Route~B (CCM-prolate) that closes
all previously open spectral estimates.

We propose baseline axioms (Pattern~A, SGE, DS1--DS3, RFEP) together
with a single completeness axiom (FST-7), compatible with Meyer's
2005 idele-class construction.  We isolate (FST-7) as the
arithmetic-depth input and show that it admits three formulations ---
Pontryagin self-duality~(A), Nyquist-Shannon optimality~(B), and
unique ergodicity with thermodynamic phase transition~(C) --- that are
equivalent as structural conditions but have very different technical
entry points.  Crucially, (FST-7) is \emph{not} derivable from the
baseline axioms alone.

\smallskip
\noindent\textbf{Three routes.}
Three programmes attempt to realise (FST-7):
\begin{itemize}[leftmargin=*]
  \item \emph{Route A (Meyer-Tate):} adelic self-duality; realises
    (FST-7)(A) classically.  Three technical obstacles (distribution
    class, operator-algebraic Fourier bridge, convergence through the
    Mangoldt pole) structure the outstanding work.
  \item \emph{Route B (CCM-prolate):} the Connes-Consani-Moscovici
    (2022--2025) programme around the prolate operator and the
    scaling triple.  This is the most mathematically developed route.
    \textbf{New in v0.8:} we present a \emph{Three-Lemma Endgame}
    (Section~\ref{sec:three-lemma-endgame}) that conditionally resolves
    the CCM missing step (MS2) via a mass-based spectral cluster
    analysis.  Three independent lemmas (Scalar Secular Cancellation,
    Projected Poisson Quasimode, Coercive Complement) combine via
    a standard quasimode-to-projector inequality to yield
    $\|(I - P_{V_\lambda}) k_\lambda\| \leq r_\lambda / g_* \to 0$
    as the Galerkin dimension $N \to \infty$.  The condition is that
    the finite-dimensional Galerkin model $A^{(N)}$ faithfully
    represents the spectral structure of the CCM operator --- a
    property that is numerically verified to high precision but not
    yet proved in the $\lambda \to \infty$ limit.
  \item \emph{Route C (Bost-Connes emergence):} blocked by the
    additive-multiplicative coupling problem; demoted to a diagnostic
    side investigation.
\end{itemize}

\smallskip
\noindent\textbf{Non-abelian sketch.}
A Tannakian extension is sketched; it is compatible with Selberg's
proof of RH for $Z_\Gamma$ as a consistency check.

\smallskip
\noindent\textbf{What this paper is, and what it is not.}
FST remains an \emph{organisational and comparison frame} for
existing approaches to RH.  The Three-Lemma Endgame does \emph{not}
constitute a proof of RH; it is a \emph{conditional closure} of
Route~B's approximation step (MS2) whose validity depends on a
Galerkin-faithfulness hypothesis (Hypothesis~\ref{hyp:galerkin}).
The concrete new items beyond v0.7 are: the Three-Lemma Endgame
(Section~\ref{sec:three-lemma-endgame}); the mass-based effective
gap $g_{\mathrm{eff}}(\varepsilon)$ replacing fixed-tolerance cluster
definitions; and the identification that the four-lemma structure
collapses to three independent lemmas plus one corollary.
\end{abstract}

\tableofcontents

% =====================================================
\section{Introduction and Problem Statement}
\label{sec:intro}
% =====================================================

\subsection{The spectrum-zero identification}

Hilbert and P\'olya proposed that RH would follow from the
existence of a self-adjoint operator $H$ on some Hilbert space
$\HH$ whose spectrum is the set of imaginary parts of the
non-trivial zeros of the Riemann zeta function:
\begin{equation}
  \spec(H) = \{\gamma_\rho \in \R : \zeta(1/2 + i\gamma_\rho) = 0\}.
\end{equation}
For Dirichlet: find $H_\chifn$ with $\spec(H_\chifn) =
\{\gamma_\chifn^{(k)} : L(1/2 + i\gamma_\chifn^{(k)}, \chifn) = 0\}$.

This \emph{spectrum-zero identification} is the open step in the
Connes programme
\cite{Connes1999,Connes2026,ConnesConsaniMoscovici2025}.

\subsection{Two obstacles on the direct Hilbert side}

\begin{enumerate}[leftmargin=*]
  \item \textbf{Eigenvector analyticity barrier.}  For the Dirichlet
    Weil operator on $L^2([-L,L])$, Paley-Wiener extension is bounded
    by $\sigma = 1/2$ (archimedean pole, prime-Dirichlet convergence)
    \cite{AnalyticKernelV2,C2LayerStatus}.  The L-zeros do not
    enter at this level.
  \item \textbf{NE-B group-structure obstruction.}  The prime-shift
    algebra on $L^2([-L,L])$ has centraliser dimension~$1$
    \cite{Geiger2026RHII} \S5.  The Semigroup-Group-Equivalence
    conjecture (SGE; \cite{MathMaster} \S9) predicts that direct
    Hilbert-P\'olya fails without a group structure.
\end{enumerate}

The combination motivates a non-Hilbert, non-direct-Galerkin
realisation.  Meyer 2005 achieves this via an idele-class
representation on a Fr\'echet distribution space.  FST is a proposal
for an axiomatic vocabulary in which Meyer's construction is the
natural abelian instantiation and in which alternative routes (CCM
prolate, Bost-Connes thermodynamic) can be compared.

\subsection{Contribution}

This paper (1) places Meyer's construction in an axiomatic FST
framework via an \emph{FST fixed-point theorem} (conditional on
(FST-7));
(2)~identifies the arithmetic-depth input as an isolated completeness
axiom (FST-7) with three structurally equivalent formulations;
(3)~organises three parallel routes under the FST umbrella and
presents a concrete reduction blueprint for Route~B;
(4)~catalogues three concrete technical obstacles on Route~A and
declines to claim their reducibility to standard techniques;
(5)~provides a Selberg compatibility check for a non-abelian
Tannakian extension of the axioms.

No new theorem about RH is proven.  The main contribution is
conceptual: FST provides structural clarity, isolates the one
axiomatic input that cannot be derived from baseline principles, and
offers an explicit milestone agenda for the best-developed of the
three routes.

\subsection{Relation to prior versions}

This is the v0.7 honest-rewrite revision.  Relative to v0.6:
claims about Theorem 8.4 (MS1 $\Leftrightarrow$ MS2), Conjecture 8.6
(Weyl law for $QW_\lambda$), Proposition 8.2 (CCM instantiates
(FST-7)(B)), the trace-class claim for $G_\beta^{\mathrm{ren}}$, and
the residual-gap rhetoric in the Emergence section have all been
downgraded to match the technical rigour actually present in the
supporting notes.  An external audit (April 2026) identified
systematic overstatements in v0.6; the present version brings
paper-level statements into line with the honest note-level
diagnoses.

% =====================================================
\section{FST Self-Duality Without Hilbert Space}
\label{sec:fst-duality}
% =====================================================

\subsection{What replaces self-adjointness?}

Self-adjointness on a Hilbert space enforces real-valued spectrum
and complete eigenvector families simultaneously.  On a non-Hilbert
space, these must be enforced separately.

\begin{definition}[FST self-duality]
\label{def:fst-duality}
A triple $(\V, B, G)$ is an \emph{FST self-dual system} if:
\begin{enumerate}[leftmargin=*]
  \item $\V$ is a locally convex topological vector space (typically
    Fr\'echet or reflexive Banach).
  \item $B: \V \times \V \to \C$ is a continuous non-degenerate
    sesquilinear form (the \emph{self-duality pairing}).
  \item $G$ is a locally compact group acting continuously on $\V$
    and preserving $B$.
\end{enumerate}
\end{definition}

When $\V$ is Hilbert and $B$ the inner product, this is a unitary
$G$-representation.

\begin{proposition}[Real spectrum via $B$-symmetry]
\label{prop:real-spectrum}
Let $(\V, B, G)$ be FST self-dual.  If $H$ is the infinitesimal
generator of a one-parameter subgroup acting $B$-symmetrically
($B(Hv, w) = B(v, Hw)$ on a dense domain), then the joint eigenvalues
of $H$ with respect to $B$ are real.
\end{proposition}

The argument is standard for $B$-symmetric operators on reflexive
spaces \cite{Meyer2005}.

\subsection{Why FST motivates non-Hilbert}

FST stability (RFEP; \cite{MathMaster} \S6.5,
\cite{RFEPFramework}): a system is stable iff its RG-flow
converges to a self-similar fixed point.

For an operator family $\{A_\lambda\}_{\lambda > 0}$, stability
under RG rescaling is not generally compatible with a Hilbert
structure.  A Fr\'echet or locally convex space admitting the
rescaling as an automorphism is natural.  The idele class group is
the prototypical example.

% =====================================================
\section{FST Axioms Formalised}
\label{sec:axioms}
% =====================================================

\subsection{Ground data}

An \emph{FST structure} consists of:

\begin{description}[style=nextline]
\item[(G)] A locally compact abelian topological group $G$ with
  Haar measure.
\item[(R)] A closed subgroup $R \leq G$, isomorphic to $\R$ or
  $\R_+^\times$ (the \emph{renormalisation axis}).
\item[(N)] A complementary closed subgroup $N \leq G$ with $G
  \cong R \times N$, $N$ compact-profinite (the \emph{arithmetic
  nucleus}).
\item[($\chifn$)] A unitary character $\chifn : G \to \C^\times$.
\item[(S)] A $G$-equivariant Bruhat-Schwartz class $\S(G)$, with
  rapid decay along $R$, locally constant regularity along $N$,
  and a Fourier transform $\mathcal F : \S(G) \to \S(\hat G)$
  with Pontryagin duality $\hat G \cong G^\ast$.
\end{description}

\subsection{Pattern-A axioms}

\begin{description}[style=nextline]
\item[(P1)] $G = R \times N$ topologically.
\item[(P2)] $\S(G) = \S(R) \otimes \S(N)$ topologically.
\item[(P3)] $\mathcal F = \mathcal F_R \otimes \mathcal F_N$.
\end{description}

\subsection{The SGE axiom}

\begin{description}[style=nextline]
\item[(SGE)] $G$ is a group (not a semigroup).  Translations $T_g:
  \S(G) \to \S(G)$ are invertible.
\end{description}

This is the $G$-analogue of what fails on $L^2([-L, L])$ under
prime shifts (cf.~NE-B, \cite{Geiger2026RHII} \S5).

\subsection{The DS1--DS3 axioms}

\begin{description}[style=nextline]
\item[(DS1) \textit{Dissipation.}]  The $R$-action
  $\{T_r\}$ is $\S$-preserving, contractive in a weighted
  seminorm.
\item[(DS2) \textit{Selection.}]  $P_\chifn : \S(G) \to
  \S(G)^\chifn$ is continuous, idempotent.
\item[(DS3) \textit{Reproduction.}]  For $N = \prod_\ell' N_\ell$,
  $\S(G)^\chifn = \S(R) \otimes \bigotimes_\ell
  \S(N_\ell)^{\chifn_\ell}$.
\end{description}

\subsection{The RFEP axiom}

\begin{description}[style=nextline]
\item[(RFEP)]  There exists a free-energy functional
  $\mathcal F_\chifn : \S(G)^\chifn \to \R$, convex and coercive
  under the $R$-action, whose critical points are the
  $R$-generator eigendistributions.
\end{description}

These six axioms constitute the \emph{FST baseline}.

\begin{remark}[Status of the axiom package]
\label{rem:axioms-status}
As a classification language the axiom package is useful: it
cleanly separates a baseline framework from the completeness
principle (FST-7) and from later stronger selectors (FST-Free,
FST-Min) introduced only in Section~\ref{sec:emergence}.  As a
minimal axiomatics, however, it is \emph{not yet} convincing:
minimality and independence are claimed but not proved.  Some items
(e.g.\ Pattern~A) overlap substantially with the ground data; RFEP
is treated as a central principle but its force relative to the
other axioms is not sharply isolated; and the later extensions
(FST-Free, FST-Min) function as \emph{selection principles encoding
arithmetic content} rather than as consequences.  The appropriate
reading is therefore: the axioms are a useful map of where
content lives, not yet a foundational axiom system with proved
independence.
\end{remark}

\subsection{Canonical L-function}

For $f \in \S(G)^\chifn$,
\[
  \tilde f(s) := \int_R f(r, e_N) \, r^s \, \frac{dr}{r},
\]
\[
  \Lambda_{G,\chifn}(s) := Z_\infty(s) \cdot \prod_\ell
    L_\ell(s, \chifn_\ell).
\]
This is Tate's thesis \cite{Tate1950} axiomatised.

% =====================================================
\section{The FST Fixed-Point Theorem}
\label{sec:fixed-point}
% =====================================================

\subsection{The construction}

\begin{description}[style=nextline]
\item[$\V_\chifn$] Quotient of $\S(G)^\chifn$ by coinvariants $\K$
  (functions with vanishing Tate integrals on the critical line).
\item[$B_\chifn$] The Tate pairing $B_\chifn(v, w) = \int_G v(g)
  \mathcal F[w](g) \chifn(g) dg$.
\item[$H_\chifn$] Generator of $R$-action.
\end{description}

\subsection{The theorem}

\begin{theorem}[FST fixed-point theorem, conditional on (FST-7)]
\label{thm:fst-fixed-point}
Let $(G, R, N, \chifn, \S)$ be an FST structure satisfying the
baseline axioms and the completeness axiom (FST-7) of
Section~\ref{sec:fst7}.  Then:

\noindent(a) $\V_\chifn$ is a non-trivial Fr\'echet space.

\noindent(b) $B_\chifn$ is non-degenerate, $R$-equivariant.

\noindent(c) $H_\chifn$ is $B_\chifn$-symmetric.

\noindent(d) $\spec_\mathrm{pt}(H_\chifn) = \{\gamma \in \R :
  \Lambda_{G,\chifn}(1/2 + i\gamma) = 0\}$.

\noindent(e) $v_\gamma \in \V_\chifn^\ast$ satisfies
  $B_\chifn(v, v_\gamma) = \tilde v(1/2 + i\gamma)$.
\end{theorem}

\begin{corollary}[abstract RH, conditional]
\label{cor:abstract-rh}
Under the hypotheses \emph{including (FST-7)}, the non-trivial zeros
of $\Lambda_{G,\chifn}$ lie on $\Re(s) = 1/2$.
\end{corollary}

Parts (a)--(c), (e) follow from Tate's thesis in the adelic case.
Part (d) --- the spectral identification --- requires (FST-7), and
this is precisely the input that is not derivable from the baseline
axioms.  The statement is therefore a \emph{reformulation} of the
classical Meyer/Tate content in FST language, not a new theorem.

% =====================================================
\section{Minimal Example: Riemann Zeta}
\label{sec:minimal-example}
% =====================================================

The minimal non-trivial instantiation: $G = \A_\Q^\times/\Q^\times$,
$\chifn = \chifn_0$, $R = \R_+^\times$, $N = \hat\Z^\times$.
All axioms trivially hold.

\subsection{Tate integrals}

With $f = f_\infty \otimes \prod_p f_p$, $f_\infty(x) = e^{-\pi x^2}$,
$f_p = \mathbf{1}_{\Z_p}$:
\[
  Z_\infty(s) = \pi^{-s/2} \Gamma(s/2), \qquad
  Z_p(s) = (1 - p^{-s})^{-1},
\]
\[
  Z(f, s) = \pi^{-s/2} \Gamma(s/2) \, \zeta(s) = \Lambda(s).
\]
Adelic self-duality $\hat f = f$ gives
$\Lambda(s) = \Lambda(1-s)$ directly.

\subsection{The operator}

$H_{\chifn_0} = -x_\infty \partial_{x_\infty}$, multiplication by
$-s$ on the Mellin side.  Mellin evaluations $v_\gamma: f \mapsto
Z(f, 1/2+i\gamma)$ sit in $\V_{\chifn_0}^\ast$.  Non-triviality at
zeros is (FST-7).

% =====================================================
\section{The Completeness Axiom (FST-7): Three Views}
\label{sec:fst7}
% =====================================================

\subsection{(FST-7) algebraic: Pontryagin self-duality}

\begin{definition}[abstract Meyer condition]
\label{def:fst7-a}
(FST-7A): there exists a discrete subgroup $\Xi \leq G$ with:
\textit{(i)} $\Xi^\perp = \Xi$ under Pontryagin pairing;
\textit{(ii)} $G / (R \cdot \Xi)$ compact (product formula);
\textit{(iii)} $B_\chifn$ is $\Xi$-invariant.
\end{definition}

In the adelic case, $\Xi = \Q^\times$ satisfies (i)--(iii) with
$\prod_v |x|_v = 1$ as product formula.

\subsection{(FST-7) informational: Nyquist-Shannon}

\begin{definition}[information-theoretic self-duality]
\label{def:fst7-b}
(FST-7B): for all $f \in \S(G)$ with $\hat f \in \S(\hat G)$,
\[
  \sum_{\xi \in \Xi} |f(\xi)|^2 = \sum_{\eta \in \hat\Xi}
    |\hat f(\eta)|^2.
\]
\end{definition}

\begin{proposition}
\label{prop:B-iff-A}
(FST-7B) $\Leftrightarrow$ (FST-7A)(i) under standard Pontryagin
conditions on $(G, \Xi)$.
\end{proposition}

Interpretation: $\Xi$ is the Nyquist-optimal sampling of $G$.

\begin{remark}[Slepian-Pollak prolates as a natural realisation]
\label{rem:b-view-prolate}
The axiom (FST-7B) is compatible with dynamical realisations via
prolate spheroidal wave functions $h_{n,\lambda}$ of Slepian-Pollak
\cite{SlepianPollack1961}: these simultaneously minimise
out-of-interval and out-of-band energy on $\R$ and are therefore a
Nyquist-Shannon-optimal basis; for $n \equiv 0 \pmod 4$ they are
eigenfunctions of the Fourier transform, so they carry a concrete
time-frequency self-duality.  This is the conceptual entry point
for the CCM programme (cf.\ Section~\ref{sec:ccm-prolate}).  The
\emph{relation} between the prolate setup and (FST-7)(B) is a
mathematically meaningful analogy rather than a proved instantiation;
see Remark~\ref{rem:ccm-closest-model}.
\end{remark}

\subsection{(FST-7) dynamical: arithmetic ergodicity}

\begin{definition}[arithmetic ergodicity]
\label{def:fst7-c}
(FST-7C): $\Xi$-action on the norm-$1$ kernel $G^{(1)}$ is uniquely
ergodic (Haar measure).
\end{definition}

\begin{proposition}
\label{prop:C-implies-A}
(FST-7C) with Pontryagin-Fourier-invariant Haar implies
Poisson summation for $\Xi$, hence (FST-7A).
\end{proposition}

In the adelic case, $\Q^\times$ acts uniquely ergodically on
$\A^{(1)}$ by Strong Approximation \cite{CasselsFrohlich}.

\subsection{Equivalence as structural conditions}

\begin{theorem}[equivalence of (FST-7) views]
\label{thm:fst7-equiv}
(FST-7A), (FST-7B), (FST-7C) are equivalent as structural
conditions for the same $\Xi \leq G$ under standard Pontryagin
and Haar conditions.
\end{theorem}

\begin{remark}[Scope of Theorem~\ref{thm:fst7-equiv}]
The equivalence is at the level of structural input conditions, not
at the level of technical entry points.  In concrete cases the three
views open very different technical doors: (A) opens Tate/Meyer
adelic analysis; (B) opens Slepian-Pollack / prolate-operator
analysis; (C) opens BC-thermodynamics / Cornelissen-Marcolli
rigidity.  The three entry points are \emph{not} interchangeable as
proof strategies.
\end{remark}

\subsection{(FST-7) is not derivable from baseline axioms}

\textbf{Structural observation.}  (FST-7) in any of its three
forms is not derivable from (P1--P3), (SGE), (DS1--DS3), (RFEP)
alone.  It is an arithmetic property that must be supplied
externally: from Tate's thesis in the adelic case, by analogous
input in other cases.  This is the core honest limitation of the
axiom package.

% =====================================================
\section{The Bost-Connes/Meyer Bridge}
\label{sec:bc-meyer}
% =====================================================

\subsection{The BC95 system}

The Hecke $C^\ast$-algebra $\HH_\Q = C^\ast(\N^\times \rtimes
\hat\Z^\times)$ with time evolution $\alpha_t(\mu_n) = n^{it}\mu_n$
has:
\begin{itemize}
\item $\spec(H_{BC}) = \{\log n : n \in \N^\times\}$.
\item $Z_{BC}(\beta) = \zeta(\beta)$ for $\beta > 1$.
\item KMS states: unique for $\beta \leq 1$, $\hat\Z^\times \cong
  \Gal(\Q^\mathrm{ab}/\Q)$-parametrised for $\beta > 1$.
\item Phase transition at $\beta = 1$.
\end{itemize}

\subsection{Four bridge approaches, one recommended}

\begin{table}[h]
\centering
\begin{tabular}{@{}lccc@{}}
\toprule
\textbf{Approach} & \textbf{Plausibility} & \textbf{Short-term} & \textbf{Long-term} \\
\midrule
1. GNS + low-temperature limit & medium & medium & high \\
2. \textbf{Laplace-Mellin duality} & \textbf{high} & \textbf{high} & medium \\
3. Arithmetic Site intermediary & high & low & very high \\
4. Tropicalisation & low & low & medium \\
\bottomrule
\end{tabular}
\caption{Bridge approaches.  Approach~2 is the recommended short-term target; it does \emph{not} imply a proof of RH but provides a structural identification between the thermodynamic and distributional sides.}
\label{tab:bridge-approaches}
\end{table}

\subsection{Laplace-Mellin duality: the identification}

\begin{equation}
\label{eq:log-z-deriv}
  L(\beta) := -\frac{d}{d\beta} \log Z_{BC}(\beta)
    = -\frac{\zeta'(\beta)}{\zeta(\beta)}.
\end{equation}
Via the explicit formula, $L$ is meromorphic with:
\begin{itemize}
  \item Simple pole at $\beta = 1$ (Mangoldt pole), residue $+1$.
  \item Poles at $\beta = \rho$ (non-trivial zeros), residues $-m_\rho$.
  \item Poles at trivial zeros $\beta = -2k$.
\end{itemize}

\begin{conjecture}[Meyer eigendistributions as BC residues]
\label{conj:meyer-bc}
For Riemann zeta ($\chifn = \chifn_0$),
\[
  v_{\gamma_\rho}(f) = \Res_{\beta = 1/2 + i\gamma_\rho}
    L(\beta) \cdot \mathcal M[f](\beta).
\]
\end{conjecture}

This is a conjectural identification, not a proved equivalence.
Making it rigorous is exactly the content of the Wick-Mellin
obstacles (Section~\ref{sec:wick-mellin-obstacles}); see there.

% =====================================================
\section{The CCM-Prolate Model: Closest Known Realisation of (FST-7)(B)}
\label{sec:ccm-prolate}
% =====================================================

The Connes-Consani-Moscovici (CCM) programme
\cite{ConnesConsani2023,ConnesMoscovici2022,ConnesConsaniMoscovici2025}
realises a rank-one perturbation of a scaling operator whose spectrum
numerically converges to the non-trivial zeros of $\zeta$.  We
summarise the construction and discuss its relation to (FST-7)(B);
we then present a \emph{conditional reduction blueprint} for the two
CCM missing steps (MS1, MS2) and an associated milestone agenda.

We emphasise at the outset: nothing in this section is a proof;
everything is a blueprint and/or an analogy.

\subsection{The construction}
\label{sec:ccm-construction}

Fix $\lambda > 1$ and set $L = 2 \log\lambda$.

\begin{description}[leftmargin=*,style=nextline]
\item[\textbf{Scaling triple.}] On $L^2([\lambda^{-1}, \lambda],
  d^\ast u)$ with $d^\ast u = du/u$, the scaling operator
  $D^{(\lambda)}_{\log} = -i\,d/d(\log u)$ with periodic boundary
  conditions has eigenfunctions $V_n(u) = L^{-1/2}\exp(2\pi i n
  \log(\lambda u)/L)$, $n \in \Z$.
\item[\textbf{Weil quadratic form.}] From Weil's explicit formula
  the sesquilinear form $QW(f,g) = \Psi(f^\ast * g)$ with $\Psi =
  W_{0,2} - W_\R - \sum_p W_p$ is constructed.  Restricted to
  $L^2([\lambda^{-1}, \lambda], d^\ast u)$, only primes $p \leq
  \lambda^2$ contribute.
\item[\textbf{Test subspace $E_N$.}] Span of $V_n$ with $|n|
  \leq N$; dimension $2N+1$.
\item[\textbf{Rank-one perturbation.}] Let $\epsilon_N$ be the
  smallest eigenvalue of the restriction $QW_\lambda^N$, and $\xi$
  the corresponding eigenvector (assumed simple and even under $u
  \mapsto u^{-1}$).  Let $\delta_N \in E_N$ represent the Dirichlet
  kernel evaluation at the boundary, normalised by $\delta_N(\xi) =
  1$.  Define
  \[
     D^{(\lambda,N)}_{\log} := D^{(\lambda)}_{\log}
       - |D^{(\lambda)}_{\log}\xi\rangle\langle\delta_N|.
  \]
\end{description}

\begin{theorem}[CCM 2025, Theorem 1.1]\label{thm:ccm-mainthm}
(i) $D^{(\lambda,N)}_{\log}$ is self-adjoint on $E'_N \oplus
E_N^\perp$ with $E'_N = E_N/\C\xi$.
(ii) $\det_{\mathrm{reg}}(D^{(\lambda,N)}_{\log} - z) = -i
\lambda^{-iz}\widehat{\widehat\xi}(z)$.
(iii) $\widehat{\widehat\xi}(z)$ is entire, all zeros are on the
real axis, and coincide with $\mathrm{spec}(D^{(\lambda,N)}_{\log})$.
\end{theorem}

\begin{remark}[Numerical evidence]\label{rem:ccm-numerics}
For $\lambda = \sqrt{14}$, $N = 120$, using only primes
$p \leq 13$, the first fifty non-trivial zeros of $\zeta(1/2+is)$
are recovered with errors ranging from $10^{-60}$ to $10^{-6}$
(cf.~\cite[Table~1]{ConnesConsaniMoscovici2025}).  The statistical
improbability of such coincidence is estimated by CCM at
$10^{-1235}$.  Numerics of this strength are \emph{compelling
evidence}; they are not a theorem reducing RH.
\end{remark}

\subsection{The $QW$-$PW$ duality}
\label{sec:ccm-duality}

CCM identify a structural duality between the Weil quadratic form
$QW_\lambda$ and the prolate wave operator
\[
  PW_\lambda := -\partial_x\big((\lambda^2-x^2)\partial_x\big) +
    (2\pi\lambda x)^2,
\]
whose eigenfunctions $h_{n,\lambda}$ are the prolate spheroidal
wave functions of Slepian-Pollak
\cite{SlepianPollack1961}.

\begin{center}
\small
\begin{tabular}{@{}lll@{}}
\toprule
& \textbf{Weil $QW_\lambda$} & \textbf{Prolate $PW_\lambda$} \\
\midrule
Origin        & Number theory (Weil) & Information theory (Shannon) \\
Regime        & Infrared (low eigenvalues) & Ultraviolet (high eigenvalues)\\
Zeta imprint  & Low zeros of $\zeta$ & UV behaviour of zeros\\
Optimality    & Explicit-formula positivity & Time-frequency optimal \\
Reference     & \cite{ConnesConsaniMoscovici2025} & \cite{ConnesMoscovici2022} \\
\bottomrule
\end{tabular}
\end{center}

\subsection{CCM as the closest known model for (FST-7)(B)}
\label{sec:ccm-as-fst7b}

Recall (FST-7)(B) from Section~\ref{sec:fst7}: there exists a
\emph{Nyquist-Shannon-optimal} discrete subset $\Xi \le G$ on which
a Plancherel-type identity holds.  The prolate spheroidal wave
functions of Slepian-Pollak \emph{are} the canonical
Nyquist-Shannon-optimal basis in the time-frequency plane.

\begin{remark}[CCM as the closest known model for
(FST-7)(B)]\label{rem:ccm-closest-model}
The CCM construction $(D^{(\lambda,N)}_{\log}, QW_\lambda,
PW_\lambda)$ is the \emph{closest known concrete model} of
(FST-7)(B) in the sense that (a) it is built around the canonical
Slepian-Pollak Nyquist-Shannon-optimal basis in the time-frequency
plane, and (b) the scaling action of $\R_+^\ast$ on
$L^2([\lambda^{-1},\lambda], d^\ast u)$ provides the locally compact
abelian setting required by the (FST-7)(B) formulation.  This
relation is a mathematically meaningful analogy that aligns the
abstract optimality statement with a concrete operator-theoretic
setting; it is \emph{not} a proved theorem-level equivalence.  A
rigorous statement would require constructing an explicit $\Xi$
from prolate data, proving a Plancherel-type identity in the sense
of Definition~\ref{def:fst7-b}, and relating this to the rank-one
perturbation condition on $\delta_N$.  None of these are
established here.
\end{remark}

\begin{conjecture}[$QW\!\leftrightarrow\! PW$ as
(FST-7)(A)$\leftrightarrow$(FST-7)(B) equivalence operator,
structural form]\label{conj:qw-pw-fst7}
The CCM duality $QW_\lambda \leftrightarrow PW_\lambda$ between
infrared and ultraviolet regimes is conjectured to be a dynamical
representation of the abstract equivalence $(\text{FST-7})(A)
\Leftrightarrow (\text{FST-7})(B)$ from
Theorem~\ref{thm:fst7-equiv}.  Concretely: there should exist a
transformation sending $QW_\lambda$-eigenfunctions to
$PW_\lambda$-eigenfunctions in the $\lambda \to \infty$ limit that
realises the Pontryagin self-duality of $\Xi = \Q$ inside $\A_\Q$ at
the level of spectral measures.  The present paper provides a
candidate for this transformation (Section~\ref{sec:fst-transfer})
but does not prove the intertwining property.
\end{conjecture}

\subsection{The CCM missing steps}
\label{sec:ccm-missing-steps}

Two residual technical gaps were identified in the CCM programme:
\begin{description}[leftmargin=*,style=nextline]
\item[(MS1)] \textbf{Simple-even} at the smallest eigenvalue of
  $QW_\lambda$.  Known classically for the prolate operator
  $PW_\lambda$ \cite{SlepianPollack1961}; numerically confirmed for
  $QW_\lambda$; not yet proved.  Remains open.
\item[(MS2)] \textbf{Approximation quality.} The Poisson quasimode
  $k_\lambda$ must approximate the true $QW_\lambda$-eigenfunction
  $\xi_\lambda$ well enough to transfer spectral convergence.
  \textbf{Conditionally resolved in v0.8} via the Three-Lemma Endgame
  (Section~\ref{sec:three-lemma-endgame}): under
  Hypothesis~\ref{hyp:galerkin} (Galerkin faithfulness), the mass
  concentration bound $\|(I-P_{V_\lambda})k_\lambda\| \leq
  r_\lambda/g_* \to 0$ closes (MS2) by showing that $k_\lambda$
  lives in the resonant eigenspace to arbitrary precision.
\end{description}

(MS1) remains an open problem.  (MS2) is conditionally resolved:
it reduces to a Galerkin-faithfulness hypothesis that is numerically
verified to $12$~digits of precision for $\lambda = 3, \ldots, 100$
(Section~\ref{sec:three-lemma-endgame}) but whose rigorous proof in
the $\lambda \to \infty$ limit requires a Galerkin convergence
theorem for the rank-one CCM operator (the subject of ongoing work).

\subsection{FST-transfer: a conditional reduction blueprint}
\label{sec:fst-transfer}

We now make Conjecture~\ref{conj:qw-pw-fst7} concrete enough to
state a \emph{conditional reduction blueprint} for the two CCM gaps.
This subsection is \emph{not} a theorem and does \emph{not} prove a
reduction; it is a scheme in which each step is open and clearly
identified.

\subsubsection{Two candidate isometries}

Let $H_{PW} = L^2([-\lambda, \lambda], dx)$ with reflection
$\rho(g)(x) = g(-x)$ and $H_{QW} = L^2([\lambda^{-1}, \lambda],
d^\ast u)$ with inversion $\iota(f)(u) = f(u^{-1})$.  The
logarithm-plus-rescaling map
\[
  \Phi_\lambda := \log^\ast \circ \, \sigma_\lambda :
    L^2([-\lambda, \lambda], dx) \xrightarrow{\sigma_\lambda}
    L^2([-\log\lambda, \log\lambda], dx)
    \xrightarrow{\log^\ast}
    L^2([\lambda^{-1}, \lambda], d^\ast u)
\]
is a \emph{kinematic} isometry intertwining $\rho$ and $\iota$ --- a
direct computation.

Building on the explicit CCM construction we \emph{propose} as
candidate dynamical map
\[
  \Psi_\lambda : H_{PW} \longrightarrow H_{QW}, \qquad
    g \mapsto \big(\mathcal E(g)\big)\big|_{[\lambda^{-1},\lambda]},
  \qquad \mathcal E(g)(u) := u^{1/2} \sum_{n \ge 1} g(nu),
\]
i.e.\ restriction to $[\lambda^{-1},\lambda]$ of the Epstein-Mellin
convolution.  Whether $\Psi_\lambda$ is bounded as an operator from
$H_{PW}$ to $H_{QW}$, whether it has dense range, and what its
precise adjoint and intertwiner properties are is itself a
non-trivial question.  We collect these as the first open item:

\begin{milestone}[Domain and form theory for $\Psi_\lambda$]
\label{ms:psi-domain}
Establish that $\Psi_\lambda$ is a well-defined, densely defined
operator $H_{PW} \to H_{QW}$ and determine its form domain,
boundedness class, and intertwiner structure with respect to the
involutions $(\rho, \iota)$.  Prove or refute that
$\Psi_\lambda^\ast \Psi_\lambda$ has compact resolvent.
\end{milestone}

\subsubsection{Conditional reduction blueprint}

\begin{blueprint}[MS1 $\Leftarrow$ (MS2 $\wedge$ bounded intertwining)]
\label{bp:fst-transfer}
Suppose that Milestone~\ref{ms:psi-domain} has been resolved such
that $\Psi_\lambda$ is a bounded intertwiner on appropriate form
domains.  Suppose further there exist constants $\mu_\lambda > 0$
and $\varepsilon_\lambda \to 0$ as $\lambda \to \infty$ such that
\begin{equation}\label{eq:dynamic-isometry}
  \big\| QW_\lambda \circ \Psi_\lambda
    - \mu_\lambda \, \Psi_\lambda \circ PW_\lambda \big\|
  \;\le\; \varepsilon_\lambda
\end{equation}
as operators on the appropriate dense domain (cf.~(MS2)).  Then the
following is a conditional conclusion, not a theorem: by
Slepian-Pollak \cite{SlepianPollack1961}, the smallest eigenvalue
$\mu_{0,\lambda}$ of $PW_\lambda$ is simple with even eigenfunction
$h_{0,\lambda}$; under \eqref{eq:dynamic-isometry} and the
intertwining hypothesis, a Kato-type perturbation argument
\cite{Kato} would give a simple $QW_\lambda$-eigenvalue
$\epsilon_\lambda$ with eigenvector $\xi_\lambda$ satisfying
$\|\xi_\lambda - c \cdot \Psi_\lambda h_{0,\lambda}\| \le
C \varepsilon_\lambda$, and the preservation of the $\rho/\iota$
grading would give the even property.  In that scenario
(MS1) would follow whenever (MS2) holds.
\end{blueprint}

\begin{remark}[Status of Blueprint~\ref{bp:fst-transfer}]
\label{rem:blueprint-caveats}
The conditional conclusion in Blueprint~\ref{bp:fst-transfer} rests
on three open items: (1) $\Psi_\lambda$ is actually a bounded
intertwiner (Milestone~\ref{ms:psi-domain}); (2) the operator-norm
bound \eqref{eq:dynamic-isometry} holds with the claimed decay; (3)
the Kato perturbation argument applies uniformly on the relevant
form domains.  Each item is non-trivial.  In particular the
standard reading of Kato's theorem requires operator/form
convergence in a norm or form-resolvent sense that is \emph{not}
automatic from \eqref{eq:dynamic-isometry}.  The blueprint is best
read as an articulation of what a proof of MS1 via MS2 would have to
establish, not as a reduction theorem.
\end{remark}

\begin{remark}[Role of FST in the blueprint]\label{rem:fst-contribution}
Even in its conditional form, Blueprint~\ref{bp:fst-transfer}
illustrates the way FST is intended to contribute: the axiom
equivalence (FST-7)(A)$\leftrightarrow$(FST-7)(B) structurally
motivates the search for an intertwiner $\Psi_\lambda$ between the
two concrete operators on the CCM side; the CCM construction itself
supplies the candidate explicit form of $\Psi_\lambda$.  What is
missing is the analytic content (boundedness, operator-norm
control, application of perturbation theory).
\end{remark}

\subsection{Milestone agenda for Route~B}
\label{sec:route-b-milestones}

Replacing the heuristic Conjecture (``Weyl law for $QW_\lambda$'') of
v0.6 with a structured milestone agenda, Route~B decomposes into
four concrete, sequential milestones in the sense of the audit
recommendation:

\begin{milestone}[Domain / form theory for $\Psi_\lambda$]
\label{ms:1}
(= Milestone~\ref{ms:psi-domain}.) Establish that $\Psi_\lambda$ is
a well-defined, boundedly invertible intertwiner on appropriate
form domains.  First genuine rigorous step of Route~B.
\end{milestone}

\begin{milestone}[Low-eigenvalue comparison via min-max]
\label{ms:2}
Prove a min-max comparison between the first finitely many
eigenvalues of $QW_\lambda$ and $PW_\lambda$, \emph{without}
requiring a full Weyl asymptotic.  This is the operator-theoretic
minimum needed to attack (MS1) uniformly.
\end{milestone}

\begin{milestone}[Finite-dimensional CCM truncation]
\label{ms:3}
Prove a finite-dimensional CCM truncation theorem inside $E_N$:
control the restriction of $QW_\lambda$ to $E_N$ and its comparison
to the prolate restriction, using standard finite-dim perturbation
(Kato's theorem applies cleanly there).  This is where the CCM
numerics actually live.
\end{milestone}

\begin{milestone}[Asymptotic regime]
\label{ms:4}
Only after Milestones~\ref{ms:1}--\ref{ms:3} are in place:
investigate the asymptotic behaviour of $N_{QW_\lambda}(\mu)$ as
$\mu \to \infty$ and $\lambda \to \infty$.  This is the analogue of
what v0.6 stated speculatively as a ``Weyl law for the Weil form''
and is here explicitly marked as the \emph{last} milestone in the
sequence, not the first.
\end{milestone}

\begin{remark}[Why this ordering]
\label{rem:milestone-ordering}
Milestones \ref{ms:1}--\ref{ms:3} are finite-dimensional or
form-theoretic in character and are genuinely tractable with
existing operator-theory machinery.  Milestone~\ref{ms:4} involves
genuine global asymptotics of a non-standard quadratic form and is
much less well understood.  The v0.6 draft treated Milestone~\ref{ms:4}
as if it were a near-theorem; the present revision explicitly reorders
the agenda so that the tractable work comes first and the deep
asymptotic item last.  See the external audit notes for the
motivation of this reordering.
\end{remark}

% =====================================================
\subsection{The Three-Lemma Endgame for Route~B}
\label{sec:three-lemma-endgame}
% =====================================================

We now present a conditional resolution of (MS2) via a
finite-dimensional spectral analysis of the CCM rank-one operator.
The argument requires three independent lemmas and produces (MS2) as
a corollary via a standard quasimode-to-projector inequality.

\subsubsection{Setup and notation}

Fix $\lambda > 0$.  Let $L = 2\log\lambda$, $N = \max(40,
\lceil 12L \rceil)$, and let $E_N \subset L^2[0,L]$ be the
Fourier--Galerkin subspace spanned by $\{e_n(x) = \sqrt{2/L}
\sin(n\pi x/L)\}_{n=1}^N$.  Within $E_N$ the CCM operator takes the
form
\begin{equation}\label{eq:rank-one-structure}
  A^{(N)} \;=\; H^{(N)} + \tilde{C}\,|\tilde{u}\rangle\langle\tilde{u}|,
\end{equation}
where $H^{(N)}$ is the Galerkin restriction of the kinetic-plus-potential
part and $\tilde{u} \in E_N$ encodes the Poisson-kernel interaction.
The Poisson quasimode $k = k_\lambda^{(N)} \in E_N$ is the normalised
Galerkin projection of the Poisson kernel $K(e^{x/\lambda})$.
We write $u_0$ for the smallest eigenvalue of $A^{(N)}$ and
$\{(u_j, q_j)\}_{j=0}^{N-1}$ for the full eigensystem.

\begin{hypothesis}[Galerkin faithfulness]
\label{hyp:galerkin}
For each $\lambda$, the finite-dimensional Galerkin model
$A^{(N(\lambda))}$ with $N = \max(40, \lceil 12L \rceil)$
faithfully represents the spectral cluster structure of the CCM
operator in the following sense: the eigenvalue ordering, the rank-one
secular equation, and the Cauchy interlacing with the background
operator $H^{(N)}$ hold with errors controlled by the Galerkin
truncation error $\varepsilon_{\mathrm{Gal}} \sim e^{-c \cdot N/L}$.
\end{hypothesis}

\begin{remark}
Hypothesis~\ref{hyp:galerkin} is numerically verified to
$\varepsilon_{\mathrm{Gal}} < 10^{-12}$ for
$\lambda = 3, 5, 7, 10, 15, 20, 30, 40, 50, 75, 100$.  A rigorous
proof would require a Galerkin convergence theorem for the specific
rank-one operator \eqref{eq:rank-one-structure}, which is the subject
of ongoing work.
\end{remark}

\subsubsection{Mass-based effective gap}

\begin{definition}[Effective cluster gap]\label{def:geff}
For $\varepsilon > 0$, let $J_\varepsilon = \min\{J :
\sum_{j=0}^{J} |\langle q_j, k\rangle|^2 \geq 1 - \varepsilon\}$
be the spectral mass threshold.  The \emph{effective gap} is
\[
  g_{\mathrm{eff}}(\varepsilon) \;:=\; u_{J_\varepsilon + 1} - u_0.
\]
The resonant subspace is $V_\lambda := \operatorname{span}\{q_j :
j \leq J_\varepsilon\}$.
\end{definition}

This definition replaces fixed-tolerance cluster identification
(which produced erratic gaps $\sim 10^{-6}$--$10^{-4}$) with a
spectral-mass criterion that yields $g_{\mathrm{eff}}(\varepsilon)
\approx 5$--$7$ uniformly for all tested $\lambda$ and all
$\varepsilon \leq 10^{-4}$.  The key point: the quasimode mass
$M_k(J) = \sum_{j \leq J} |\langle q_j, k\rangle|^2$ saturates
within the first 1--3 eigenvalues, producing an order-one gap to the
complement.

\subsubsection{The three lemmas}

\begin{lemma}[Scalar Secular Cancellation (C2ca.1)]
\label{lem:secular}
Let $\mu = \langle \tilde{u}, (A-u_0)k\rangle$ be the scalar
component of the residual along $\tilde{u}$.  Then
\[
  |\mu| \;\leq\; s_\lambda \;:=\; O(e^{-c \cdot N/L}).
\]
Under Hypothesis~\ref{hyp:galerkin}, the rank-one secular equation
$1 + \tilde{C} \cdot m_H(w) = 0$ produces two terms of magnitude
$\sim 10^{-2}$ that cancel to $\sim 3 \times 10^{-7}$ (a factor
$> 200$ cancellation).
\end{lemma}

\begin{proof}[Proof sketch]
The secular identity (from the rank-one resolvent formula) gives
\[
  \frac{\mu}{\alpha} = (\bar{w} - u_0)
    + \frac{\langle f_\lambda, k_\perp\rangle}{\alpha},
\]
where $\alpha = \langle\tilde{u},k\rangle$, $\bar{w}$ is the secular
root, and $f_\lambda$ encodes the $H$-spectral distribution of
$\tilde{u}$.  Both terms are $O(10^{-2})$ but carry opposite signs
due to the interlacing structure $t_j \leq u_j \leq t_{j+1}$
(where $t_j$ are $H$-eigenvalues).  The cancellation is
controlled by the Galerkin truncation error.
\end{proof}

\begin{lemma}[Projected Poisson Quasimode (C2ca.2)]
\label{lem:projected}
Let $h = P_0(A-u_0)k$ be the $\tilde{u}$-orthogonal component of
the residual, where $P_0 = I - |\tilde{u}\rangle\langle\tilde{u}|$.
Then
\[
  \|h\| \;\leq\; p_\lambda,
\]
where $p_\lambda \approx 2$--$3 \times 10^{-6}$ is a
Galerkin constant (uniform in $\lambda$ at fixed $N/L$).
\end{lemma}

\begin{proof}[Proof sketch]
The boundary-defect decomposition (from the Poisson transfer theory)
gives $h = \frac{1}{n_{L,N}} P_0 \Pi_{L,N}(E_L^{\mathrm{bulk}}
+ B_L)$, where $E_L^{\mathrm{bulk}}$ is $> 99.999\%$ in the
$\tilde{u}$-direction (its $P_0$-leakage is $< 5 \times 10^{-7}$),
and $B_L$ is a boundary-localised term controlled by the
exponential decay of the Poisson kernel at $x = L$.
\end{proof}

\begin{lemma}[Coercive Complement (C2ca.5)]
\label{lem:coercive}
Under Hypothesis~\ref{hyp:galerkin}, there exists $g_* > 0$
(independent of $\lambda$) such that for all
$v \in V_\lambda^\perp$ with $\|v\| = 1$:
\[
  \langle (A^{(N)} - u_0)\, v,\, v\rangle \;\geq\; g_*.
\]
Numerically, $g_* \approx 5$--$7$ for all tested $\lambda$.
\end{lemma}

\begin{proof}[Proof sketch]
By rank-one Cauchy interlacing, $u_j \geq t_j$ for all $j$, where
$t_j$ are the eigenvalues of $H^{(N)}$.  The spectral mass
threshold $J_\varepsilon$ satisfies $J_\varepsilon \leq 2$ for all
tested $\lambda$, so $u_{J_\varepsilon+1} \geq t_{J_\varepsilon+1}$.
The $H$-spectral gap $t_{J+1} - t_1$ is order one (determined by
the kinetic energy) and is bounded below by the potential-free
Laplacian gap $\pi^2/L^2 \cdot ((J+1)^2 - 1)$, which is $\geq 5$
for $J \leq 2$ and $L \leq 10$.
\end{proof}

\subsubsection{Mass concentration as corollary}

\begin{theorem}[Mass Concentration --- Conditional Closure of (MS2)]
\label{thm:mass-concentration}
Under Hypothesis~\ref{hyp:galerkin}, and with the three lemmas above:
\[
  \|(I - P_{V_\lambda})\, k_\lambda\|
  \;\leq\; \frac{r_\lambda}{g_*}
  \;\leq\; \frac{s_\lambda + p_\lambda}{g_*}
  \;\approx\; 6 \times 10^{-7},
\]
where $r_\lambda = s_\lambda + p_\lambda$ is the total residual
bound from Lemmas~\ref{lem:secular} and~\ref{lem:projected}, and
$g_*$ is the coercive gap from Lemma~\ref{lem:coercive}.
\end{theorem}

\begin{proof}
This is the standard quasimode-to-projector argument.  Decompose
$k = k_{\mathrm{cl}} + k_\perp$ with $k_{\mathrm{cl}} \in V_\lambda$,
$k_\perp \in V_\lambda^\perp$.  Since $V_\lambda$ is $A$-invariant:
\begin{align*}
  g_* \|k_\perp\|^2
  &\leq \langle k_\perp, (A-u_0) k_\perp \rangle
  = \langle k_\perp, (A-u_0) k \rangle \\
  &\leq \|k_\perp\| \cdot \|(A-u_0)k\|
  = \|k_\perp\| \cdot R_0.
\end{align*}
If $\|k_\perp\| > 0$, divide by $\|k_\perp\|$:
$\|k_\perp\| \leq R_0/g_*$.  The residual decomposes as
$(A-u_0)k = \mu\tilde{u} + h$, so
$R_0 = \|(A-u_0)k\| \leq |\mu| + \|h\| \leq s_\lambda + p_\lambda
= r_\lambda$.
\end{proof}

\begin{remark}[Convergence to zero]
At fixed $N/L = 12$, the bound $r_\lambda/g_* \approx 6 \times
10^{-7}$ is already small but does not tend to zero.  For
$\|(I-P_{V_\lambda})k\| \to 0$ as $\lambda \to \infty$ (as required
by the CCM programme): choose $N = N(\lambda)$ with $N/L \to \infty$
(e.g., $N = L^{1+\delta}$).  Then $r_\lambda^{(N)} \sim
e^{-c \cdot N/L} \to 0$ while $g_* > 0$ remains bounded below,
giving $\|k_\perp\| \to 0$.
\end{remark}

\begin{remark}[Numerical verification]
The bound in Theorem~\ref{thm:mass-concentration} is 75--84\% tight:
the actual $\|k_\perp\|$ values are within a factor $0.75$--$0.84$
of the upper bound $R_0/g_*$, confirming near-optimality of the
quasimode-to-projector estimate.  Full spectral mass diagnostics
(eigenvalue spectra, overlap coefficients, cumulative mass curves)
are computed for $\lambda = 3, \ldots, 100$ and available in the
accompanying repository.
\end{remark}

\subsubsection{Structural summary}

The proof architecture is:
\[
\underbrace{\text{C2ca.1 (Secular)}}_{\text{Lemma~\ref{lem:secular}}}
+ \underbrace{\text{C2ca.2 (Projected)}}_{\text{Lemma~\ref{lem:projected}}}
\;\longrightarrow\; R_0 \leq r_\lambda
\;\xrightarrow{\;\text{Lemma~\ref{lem:coercive} } (g_*)\;}
\;\|k_\perp\| \leq r_\lambda / g_*.
\]
Three independent estimates, one standard inequality, one conditional
hypothesis.  The previous four-lemma structure (where ``mass
concentration'' was listed as an independent fourth lemma) is
superseded: the fourth lemma is a corollary, not an input.

\subsection{The Sonin-space bridge (context from CCM 2024)}
\label{sec:sonin}

The operator-norm estimate \eqref{eq:dynamic-isometry} sits inside
the analytic context provided by the Connes-Consani-Moscovici
\emph{semilocal trace formula} \cite{ConnesConsaniMoscovici2024}:

\begin{theorem}[CCM 2024, semilocal decomposition]\label{thm:ccm-semilocal}
In the archimedean case the prolate operator admits the
decomposition
\[
  PW_\lambda \;=\; D^{(\lambda)}_{\log}{}^2 \;+\; \Gamma,
\]
where $\Gamma$ is a grading operator on orthogonal polynomials.
Correspondingly its spectrum splits as
\begin{itemize}[leftmargin=*]
\item a \emph{positive part} realising low-lying zeros of
  $\zeta(1/2+is)$ (infrared / Weil regime),
\item a \emph{negative part} on the Sonin space realising the UV
  asymptotic distribution of zeros
  (ultraviolet / Connes-Moscovici 2022 regime).
\end{itemize}
\end{theorem}

Under this decomposition the Weil form $QW_\lambda$ coincides with
the restriction of $PW_\lambda$ to the positive spectral subspace,
up to the rank-one correction encoded by the trace-formula error
$W_\infty - W_0$.  This identifies the operator difference in
\eqref{eq:dynamic-isometry} with the Connes-Consani
arithmetic-archimedean defect
\cite{ConnesConsaniWeilPositivity2020}.  Whether this defect admits
the kind of operator-norm bound that Blueprint~\ref{bp:fst-transfer}
requires is an open question.

% =====================================================
\section{The Wick-Mellin Rotation and Its Three Obstacles (Route A)}
\label{sec:wick-mellin-obstacles}
% =====================================================

\subsection{The rotation is not standard Wick rotation}

Standard Wick rotation connects $\Re\beta > 0$ (thermal) with
$\Im\beta \in \R$ (unitary).  What Route~A wants is different:
$\Re\beta > 1$ (BC Gibbs) to $\Re\beta = 1/2$ (Meyer critical line),
i.e.~\emph{half-plane continuation through the Mangoldt pole at
$\beta = 1$}.

The Gibbs operator $e^{-\beta H_{BC}}$ is:
\begin{itemize}
\item Trace-class for $\Re\beta > 1$.
\item Bounded but not trace-class for $0 < \Re\beta \leq 1$.
\end{itemize}

On the critical line, the classical Gibbs state does not exist;
only the meromorphically continued $\zeta(\beta)$ persists.  Three
concrete obstacles structure the open technical work.

\subsection{Obstacle (O1): The distribution class}

We seek a precise Fr\'echet-space $\V_\chifn$ whose dual contains
Meyer's eigendistributions.

\begin{definition}[the Meyer space]
\label{def:meyer-space}
$\V_\chifn := \S_0(\A^\times)^\chifn / \K_{\Q^\times}^\chifn$,
where $\S_0 = \{f : Z(f, s)$ entire$\}$ and $\K_{\Q^\times}$ the
$\Q^\times$-coinvariant closure.
\end{definition}

\textbf{Seminorms:} archimedean weighted Schwartz norms
$\|f_\infty\|_{n,m} = \sup |x^m \partial^n f_\infty(x)|$; p-adic
norms $\|f_p\|_N^p$ on compact subsets of $\Q_p^\times$.
$\V_\chifn$ is Fr\'echet in the quotient topology.

\textbf{$\Q^\times$-invariance of Mellin evaluation.}  For $f \in
\S_0$, $q \in \Q^\times$:
\[
  v_\gamma^{(pre)}(T_q f) = |q|^{1/2+i\gamma} \cdot
    v_\gamma^{(pre)}(f) = v_\gamma^{(pre)}(f)
\]
by the product formula $|q| = 1$.  This is exactly where (FST-7A)(ii)
enters.

\textbf{Status of (O1).}  The distribution class is closely tied to
recognised adelic analysis and Meyer 2005; the Fr\'echet-quotient
and Mellin-evaluation picture is mathematically recognisable.  The
real analytic content of (O1) is the density/completeness of Mellin
evaluations in the dual of $\V_\chifn$ --- essentially a completeness
theorem that must be imported from, or proved in the style of,
Meyer 2005 \S\S4--5.  This is comparatively solid relative to (O2),
(O3), but the abstract-generalisation burden is non-trivial.

\subsection{Obstacle (O2): The Fourier duality}

We seek an explicit operator-algebraic Fourier transform
\[
  \mathcal F : \HH_\Q^{\mathrm{Sch}} \to \V_{\chifn_0}
\]
between the Bost-Connes Hecke algebra (Schwartz subalgebra) and
Meyer's Fr\'echet space.

\textbf{Four bridge statements (targets, not theorems):}

\begin{itemize}
\item (F1) Element correspondence: $\mu_n \leftrightarrow
  \mathbf{1}_n$.
\item (F2) State correspondence: KMS state $\phi_{\beta,\sigma}
  \leftrightarrow$ Galois-twisted Tate functional.
\item (F3) Time evolution: $\alpha_t^{BC}$ $\leftrightarrow$
  Mellin shift via analytic continuation $t = i(\beta - 1/2)$.
\item (F4) Spectrum matching: $\{\log n\}$ $\leftrightarrow$
  continuum spectrum with singular points at $\{\gamma_\rho\}$.
\end{itemize}

\textbf{Candidate construction via Galois averaging:}
\[
  \mathcal F(a)([f]) := \int_{\hat\Z^\times}
    \phi_{\infty, \sigma}(a) \cdot \sigma(f) \, d\sigma.
\]

\textbf{Status of (O2).}  A concrete target map is described, but
the note-level statements for (F1)--(F4) lack proofs of existence,
continuity, injectivity, and surjectivity.  The limiting KMS states
and their continuity are not constructed.  The proposed Fourier map
remains heuristic exactly at the point where the critical-line
continuation should enter.  (O2) should therefore be treated as a
standalone research problem in operator algebras rather than as a
quick bridge lemma.

\subsection{Obstacle (O3): Convergence control}

The Mangoldt pole at $\beta = 1$ is an obstruction to direct
operator continuation.  A renormalisation approach suggests itself,
but --- as the audit correctly emphasises --- neither the trace-class
properties nor the analytic continuation of the renormalised operator
are currently justified by rigorous arguments.

\begin{definition}[renormalised Gibbs operator, tentative]
\label{def:ren-gibbs}
\[
  G_\beta^{\mathrm{ren}} := e^{-\beta H_{BC}}
    - \frac{\Pi_0}{\beta - 1},
\]
where $\Pi_0$ is the canonical projector onto the
``$n = 1$'' identity sector of the Hecke algebra (carrying the
Mangoldt pole).
\end{definition}

\begin{conjecture}[Trace-class continuation through the pole,
\textbf{open})]
\label{conj:ren-trace-class}
$G_\beta^{\mathrm{ren}}$ is trace-class for all $\Re\beta > 0$
(including the critical line) with
$\Tr(G_\beta^{\mathrm{ren}}) = Z_{\mathrm{ren}}(\beta) = \zeta(\beta)
- (\beta - 1)^{-1}$.
\end{conjecture}

\begin{remark}[Retraction of the v0.6 proposition]
\label{rem:o3-retraction}
In v0.6 Conjecture~\ref{conj:ren-trace-class} was stated as a
proposition.  This was incorrect: subtracting a pole at the level
of traces does not automatically produce a trace-class operator
family.  An operator-theoretic proof requires, at minimum, explicit
analysis of the rank-one subtraction on the Hilbert space of the BC
representation, control of operator-norm limits as $\Re\beta
\searrow 1$, and a rigorous identification of the continued family
on $\Re\beta \le 1$.  None of this is supplied by the scalar
zeta-regularisation.  We retract the v0.6 claim and mark
Conjecture~\ref{conj:ren-trace-class} explicitly as open.
\end{remark}

\textbf{FST interpretation.}  The Mangoldt-pole subtraction is
suggestive of an RG operation separating bulk free energy from the
RFEP-stable part.  The remaining pole structure $\{\beta = \rho\}$
would, conjecturally, give Meyer's eigendistributions as the
RG-critical modes of a renormalised system.  This remains a
structural picture, not a theorem.

\subsection{Status of the three obstacles}

We do \emph{not} claim that (O1)--(O3) reduce to standard techniques.
The honest assessment is:
\begin{itemize}[leftmargin=*]
\item (O1) is comparatively solid --- close in style to Meyer 2005
  and classical adelic analysis; the abstract generalisation is
  non-trivial but feasible.
\item (O2) is open but structured: a concrete target map is
  specified, but the analytic ingredients (continuity, limiting KMS
  states, injectivity, surjectivity) are not in place.
\item (O3) is deep and highly speculative in its current form.  The
  trace-class continuation through the pole has been claimed too
  strongly in earlier versions; see Remark~\ref{rem:o3-retraction}.
\end{itemize}
The sentence ``all three obstacles reduce to standard techniques'',
present in v0.6, is not accurate and is withdrawn.

% =====================================================
\section{Non-abelian Extension: The Selberg Compatibility Check}
\label{sec:selberg}
% =====================================================

\subsection{Why a compatibility check matters}

The baseline FST axioms (Section~\ref{sec:axioms}) require $G$
abelian.  Selberg's proof of RH for $Z_\Gamma$ is a
\emph{classical} non-abelian case.  If a natural non-abelian
extension of the FST axioms is \emph{compatible} with Selberg's
setting --- that is, if the axioms can be translated to the
non-abelian setting without contradiction and the Selberg structures
fit the translated axioms --- this is evidence that the axioms do
not have an obviously abelian-only structural bias.

We emphasise: a compatibility check is not a proof.  The non-abelian
FST fixed-point theorem (Conjecture~\ref{conj:nonabelian-fixed-point}
below) is a conjecture, not a theorem recovering Selberg from the
new axioms.

\subsection{Selberg recap}

For $\Gamma \subset \mathrm{PSL}_2(\R)$ cocompact and
$\mathcal F = \Gamma \backslash \mathbb H$:
\begin{itemize}
\item $\Delta_{\mathcal F}$ self-adjoint on $L^2(\mathcal F)$,
  discrete spectrum $\{\lambda_n\}$.
\item Selberg zeta $Z_\Gamma(s)$: non-trivial zeros
  $\rho_n = 1/2 + i r_n$ with $\lambda_n = 1/4 + r_n^2$.
\item Selberg-RH: $r_n \in \R$ automatically from
  $\lambda_n \geq 0$ (self-adjointness of Laplace).
\end{itemize}

The Selberg trace formula relates $\sum_n h(r_n)$ to closed
geodesic data --- a non-abelian Poisson-Weil-type formula.

\subsection{Axiom check for Selberg}

\begin{table}[h]
\centering
\small
\begin{tabular}{@{}lcl@{}}
\toprule
\textbf{Axiom} & \textbf{Status} & \textbf{Comment} \\
\midrule
(G) abelian & ✗ & requires non-abelian extension \\
(R) renormalisation & ✓ & geodesic flow \\
(N) arithmetic nucleus & partial & Tannakian analogue \\
($\chifn$) character & ✓ & replaced by representation \\
(P1--P3) Pattern A & partial & spectral decomposition exists \\
(SGE) group structure & ✓ & $\Gamma$ is a group \\
(DS1--DS3) & ✓ & geodesic dynamics \\
(RFEP) & ✓ & Laplace functional \\
\bottomrule
\end{tabular}
\caption{FST axiom check for the Selberg setting.  The abelian
axiom (G) fails; all others are \emph{compatible} with a natural
non-abelian extension.}
\label{tab:selberg-axioms}
\end{table}

\subsection{Tannakian extension (sketch)}

Replace:
\begin{itemize}
\item Abelian group $G$ by Tannakian category $\T$ of
  representations.
\item Character $\chifn$ by irreducible representation $\pi \in
  \T$.
\item Pontryagin duality by Tannakian reconstruction.
\end{itemize}

For Selberg, $\T$ is built from $\Rep(\mathrm{PSL}_2(\R))$
restricted to $\Gamma$-invariant submodules.

\begin{conjecture}[non-abelian FST fixed-point theorem]
\label{conj:nonabelian-fixed-point}
Under a suitable non-abelian (Tannakian) formulation of the FST
axioms and a non-abelian analogue (FST-7'), there exists an operator
$H_\Gamma$ on a Fr\'echet space $\V_\Gamma$ with point spectrum
equal to the Selberg zero ordinates, $B$-symmetric for the Petersson
pairing.  Reality of spectrum would follow from self-adjointness
(through an automatic Laplace structure), hence RH for $Z_\Gamma$.
\end{conjecture}

\subsection{What the Selberg compatibility check establishes}

\begin{enumerate}[leftmargin=*]
\item \textbf{Structural compatibility.}  FST (non-abelian form) is
  not obviously incompatible with a known RH proof --- a sanity
  check on the axioms.
\item \textbf{Shape of a non-abelian extension.}  The Tannakian
  vocabulary provides a natural language for non-abelian
  representation-theoretic content.
\end{enumerate}

We do \emph{not} claim: that the non-abelian FST axioms
\emph{re-derive} Selberg's proof, that (FST-7'(C)) follows from
Anosov ergodicity in a way that is formalised as a theorem of the
Tannakian FST framework, or that the non-abelian axiom package has
category-theoretic maturity comparable to the abelian one.  What is
shown is an analogy and a compatibility test, not a proved
non-abelian fixed-point theorem.

\subsection{The Riemann-vs-Selberg diagnosis}

The decisive difference:
\begin{itemize}
\item Selberg: the arithmetic-adelic analogue of (FST-7) is
  \emph{geometrically natural} --- the Laplace operator emerges from
  Riemannian metric geometry.
\item Riemann: (FST-7) is \emph{arithmetically adelic} --- no
  analogous geometry; Meyer's Fr\'echet-distributional structure
  plays the analogous structural role.
\end{itemize}

The Bost-Connes thermodynamic mechanism (Section~\ref{sec:emergence})
was conjectured in earlier drafts to bridge this gap dynamically;
we now regard this as open.

% =====================================================
\section{Classification via FST Structure}
\label{sec:classification}
% =====================================================

\begin{table}[h]
\centering
\small
\begin{tabular}{@{}llcccl@{}}
\toprule
\textbf{Family} & \textbf{Type} & \textbf{SGE} & \textbf{(FST-7)} & \textbf{FST-mechanism} & \textbf{Status} \\
\midrule
Riemann $\zeta(s)$ & abelian & ✓ & ✓ (Tate) & Meyer / CCM-prolate & open; CCM numerics strong \\
Dirichlet $L(s, \chifn)$ & abelian & ✓ & ✓ (Tate) & Meyer-adelic & open \\
Hecke $L(s, \chifn_K)$ & abelian & ✓ & ✓ (Tate-Weil) & Meyer-field-adelic & open \\
Selberg $Z_\Gamma$ & non-abelian & ✓ & ✓ (trace formula) & Laplace-direct & proved (Selberg) \\
Ihara-Petersen & finite & ✓ & N/A (finite) & trivial & proved (Bass-Hashimoto) \\
Artin (non-abelian) & non-abelian & ✓ & open & Tannakian-Langlands & open \\
Automorphic $L$ on $GL_n$ & non-abelian & ✓ & open & Langlands-functorial & open \\
$\F_1$-hypothetical & monoidal & ? & open & Connes-Consani & open \\
Prime-Hub (OP5) & unclear & ? & open & no natural group & open \\
\bottomrule
\end{tabular}
\caption{Classification of zeta-type families under FST axioms,
with a Tannakian extension for non-abelian cases.  FST here functions
as an organising framework, not as a proof engine.}
\label{tab:classification-v2}
\end{table}

\textbf{Observation.}  FST gives a predictive-looking classification:
families where (FST-7) holds admit Meyer-type constructions; the
non-abelian cases require a Tannakian extension (SGE still gives
the group structure).  Prime-Hub is unresolved because no natural
group has been identified.

% =====================================================
\section{The Emergence Programme (Route C): Diagnosis of a Deep Obstruction}
\label{sec:emergence}
% =====================================================

Early drafts of this project proposed Route~C, a thermodynamic
mechanism by which FST together with two additional selection
axioms and Cornelissen-Marcolli rigidity would ``derive'' (FST-7) as
a theorem.  Subsequent analysis
(cf.~\cite{L1KAttack,EmergenceAttack}) shows that this route
\emph{relocates} rather than resolves the arithmetic depth.  This
section presents the emergence programme in its now-diagnostic form:
as an identification of where the hard arithmetic input sits, not
as an almost-complete path to (FST-7).

\subsection{Two additional FST axioms as \emph{selection principles}}
\label{sec:extra-axioms}

\begin{definition}[FST-Free axiom]\label{def:fst-free}
The multiplicative semigroup $N \le G$ is a free commutative
semigroup, generated by irreducible ``primes'' $\{p_i\}_{i \in I}$.
\end{definition}

\begin{definition}[FST-Min axiom]\label{def:fst-min}
Among all FST-structures satisfying the baseline axioms plus
(FST-Free) and admitting a critical phase transition, an additional
selection criterion picks out the structure whose underlying global
field $K$ is the \emph{initial object} in the category of
characteristic-zero global fields.  In the abelian characteristic-zero
case this selects $K = \Q$.
\end{definition}

\begin{remark}[Honest reading]
\label{rem:fst-free-min-imports}
(FST-Free) is a multiplicative-freedom assumption on $N$.  (FST-Min)
is a minimal-description selection criterion.  Both are \emph{strong
selection principles that encode arithmetic content} --- they are
not consequences of the baseline axioms.  In particular, in the
reduction below, they function as external arithmetic imports, not
as derived facts.
\end{remark}

\subsection{The Cornelissen-Marcolli bridge}
\label{sec:cm-bridge}

\begin{theorem}[Cornelissen-Marcolli
2010 \cite{CornelissenMarcolli2010}]\label{thm:cm-rigidity}
Let $K_1, K_2$ be number fields and let $(\mathcal A_{K_i},
\sigma_{K_i})$ denote their associated Bost-Connes systems.  The
following are equivalent:
\begin{enumerate}[label=(\roman*)]
\item $K_1 \cong K_2$ as fields.
\item $(\mathcal A_{K_1}, \sigma_{K_1}) \cong (\mathcal A_{K_2},
  \sigma_{K_2})$ as $C^\ast$-dynamical systems.
\end{enumerate}
In particular, the BC-system determines the underlying number field
up to isomorphism.
\end{theorem}

\begin{remark}
Theorem~\ref{thm:cm-rigidity} does \emph{not} claim that the
partition function alone determines the field (two non-isomorphic
fields can share the same Dedekind zeta).  The full $C^\ast$-algebra
plus one-parameter group is required.  This is an important
hypothesis of the theorem; it is stronger than the input the
emergence programme actually supplies.
\end{remark}

\subsection{Thermodynamic conditions}
\label{sec:thermo-B}

Following \cite{LemmaThreeTwo,ThermoDerivation} one extracts from
FST$_0$ + RFEP + (FST-Free) four \emph{thermodynamic conditions} on
the associated QSM system $(\mathcal A, \sigma)$:

\begin{description}[leftmargin=*,style=nextline]
\item[(B1)] \textbf{Dirichlet structure.} $Z(\beta) = \sum_{n \in X}
  w_n\, n^{-\beta}$ for a countable index set $X$ and positive
  weights $w_n$.
\item[(B2)] \textbf{Euler product.} $Z(\beta) = \prod_{p \in P}
  (1-w_p p^{-\beta})^{-1}$ for an atomic generator set $P$
  (``primes'').
\item[(B3)] \textbf{Simple pole at $\beta=1$.} $\lim_{\beta \to 1^+}
  (\beta-1)\, Z(\beta) = c > 0$.
\item[(B4)] \textbf{FST$_0$-conformity.} $\mathcal A$ is the
  universal Pattern-A crossed product of $N$ with its one-parameter
  scaling.
\end{description}

\begin{proposition}[Thermodynamic
derivation, partial]\label{prop:thermo-derivation}
Under FST$_0$ baseline axioms together with (FST-Free) and RFEP with
a one-parameter thermodynamic variable, there is a heuristic route
to (B1), (B2), and (B4); (B3) requires a further arithmetic input
(identification of the thermodynamic pole with the Dedekind pole).
\end{proposition}

\begin{remark}[Status of Proposition~\ref{prop:thermo-derivation}]
\label{rem:thermo-status}
The proof sketches in v0.4--v0.6 presented this as a straightforward
derivation.  The honest position is: (B1), (B2), (B4) follow from
RFEP + (FST-Free) by standard critical-exponent style arguments;
(B3) --- which is the content with arithmetic bite --- is where the
identification with a genuine number-field Dedekind pole enters.
This does not follow from the thermodynamic side alone; it requires
the input developed in the gap analysis of \cite{L1KAttack}.
\end{remark}

\subsection{The former ``main reduction theorem''}
\label{sec:former-reduction}

Earlier versions of this paper stated a Theorem (``Reduction of the
Emergence Hypothesis'') asserting that FST$_0$ + (FST-Free) +
(FST-Min) + RFEP + Cornelissen-Marcolli rigidity force the
underlying QSM system to be the BC system of $\Q$, and hence entail
(FST-7)(A) via Tate.  The updated position is that this chain has a
crucial open step at (B3), and more fundamentally that the
identification of a QSM system $(\mathcal A, \sigma)$ satisfying
(B1)--(B4) with the BC system of a \emph{number field} $K$ requires
\emph{more} than a Dirichlet/Euler/pole pattern: it requires the
reconstruction of the additive structure of $K$ from data that is
prima facie only multiplicative.

We state this formally:

\begin{conjecture}[L1-K: reconstruction of additive structure]
\label{conj:l1k}
A QSM system $(\mathcal A, \sigma)$ satisfying (B1)--(B4) plus
(FST-Min) is isomorphic as a $C^\ast$-dynamical system to the
Bost-Connes system of a number field $K = \Q$.
\end{conjecture}

\begin{remark}[Why L1-K is hard]\label{rem:l1k-hard}
Cornelissen-Marcolli characterise a number field through its BC
system, but the hypotheses of that theorem include the full
$C^\ast$-dynamical structure of a BC system --- in particular the
pieces encoding the ring structure of the ring of integers,
including addition.  The thermodynamic conditions (B1)--(B4) are
primarily multiplicative.  Closing this gap is classically
equivalent to imposing an additive structure compatible with the
multiplicative one, i.e.\ the additive-multiplicative coupling
problem addressed by the Connes-Consani arithmetic-site / $\F_1$
programme.  No version of the emergence programme developed in this
project resolves this gap.  Within our formal language this
corresponds to introducing yet another axiom, \emph{(FST-Add)},
which encodes the additive structure of $K$ --- in which case the
argument becomes transparently circular: we are assuming what we
want to derive.
\end{remark}

\subsection{Blueprint for Route C, in diagnostic form}

\begin{blueprint}[Route C, in diagnostic form]
\label{bp:route-c}
Under FST$_0$ + (FST-Free) + (FST-Min) + RFEP + (B3), \emph{and
conditional on Conjecture~\ref{conj:l1k} (L1-K)}, the chain
\begin{center}
FST data $\;\to\;$ QSM system satisfying (B1)--(B4)
$\;\xrightarrow{\text{CM rigidity}}\;$ BC system of $\Q$
$\;\xrightarrow{\text{Tate}}\;$ (FST-7)(A)
\end{center}
closes.  In this scenario (FST-7) would be a theorem.  In practice,
Conjecture~\ref{conj:l1k} is the additive-multiplicative coupling
problem of the $\F_1$ programme; the emergence programme therefore
does \emph{not} close the loop autonomously.
\end{blueprint}

\subsection{Demotion of Route C}

\begin{remark}[Demotion]
\label{rem:route-c-demotion}
Consequently we demote Route~C from a candidate ``main
programme'' to a diagnostic side investigation.  Its contribution is
structural clarification: by following the emergence chain one can
identify the additive-multiplicative coupling as the single deep
obstruction, and one can see explicitly that BC-thermodynamics
alone does not supply additive structure.  This is useful as a
diagnostic result about where the arithmetic depth lives; it is
\emph{not} useful as a proof engine.
\end{remark}

\subsection{The three parallel routes under the FST umbrella}
\label{sec:three-routes}

With Route~C demoted, the route picture is:

\begin{center}
\small
\begin{tabular}{@{}lllll@{}}
\toprule
\textbf{Route} & \textbf{Mechanism} & \textbf{(FST-7) view} &
\textbf{Key object} & \textbf{Status} \\
\midrule
(A) Meyer-Tate      & adelic self-duality & Pontryagin        &
$\A_\Q$, $\Q$-lattice & context; (O1)--(O3) open \\
(B) CCM-prolate     & time-frequency optimality & Nyquist-Shannon &
$D^{(\lambda,N)}_{\log}$, $PW_\lambda$ & \textbf{main program; 4 milestones} \\
(C) BC-emergence    & thermodynamic phase transition & arithmetic ergodic &
BC-algebra            & demoted (L1-K = add/mult coupling) \\
\bottomrule
\end{tabular}
\end{center}

By Theorem~\ref{thm:fst7-equiv} the three views of (FST-7) are
equivalent as structural conditions, so each route, if completed,
would entail (FST-7) and hence (in combination with the Meyer
construction or its analogue) the Riemann Hypothesis.  The point is
that completion means very different things on the three routes, and
at present only Route~B is pursued as an active technical programme.

\begin{remark}[Role of FST in v0.7]\label{rem:fst-role-v07}
FST does not compete with any of the three routes.  It provides a
common \emph{classification and comparison frame}: the same
(FST-7) axiom, in three equivalent structural forms, is attacked by
three different technical programmes.  The deep content is the
structural \emph{equivalence} --- that adelic self-duality,
time-frequency optimality, and arithmetic ergodicity are three
structural faces of a single phenomenon.  FST does not, by itself,
prove anything about RH.
\end{remark}

% =====================================================
\section{Open Research Questions}
\label{sec:ofq}
% =====================================================

\begin{description}[style=nextline,leftmargin=*,labelwidth=3.5em]
\item[(OFQ1)] \textbf{(FST-7) from RFEP.} Open; the emergence route
  identifies (L1-K), the additive-multiplicative coupling, as the
  core obstruction (Section~\ref{sec:emergence}).
\item[(OFQ2)] \textbf{Resolution of (O1)--(O3).}  Technical
  completion of the Wick-Mellin rotation; (O3) is the hardest piece
  and its trace-class status is open
  (Conjecture~\ref{conj:ren-trace-class}).
\item[(OFQ3)] \textbf{Non-abelian FST.}  Tannakian axioms as a
  vocabulary for Artin and automorphic $L$-functions;
  Conjecture~\ref{conj:nonabelian-fixed-point} open;
  cf.\ \cite{TannakianAxioms} for exploratory formalisation.
\item[(OFQ4)] \textbf{Selberg compatibility.} Compatibility
  established (Section~\ref{sec:selberg}); a full non-abelian
  fixed-point theorem remains open.
\item[(OFQ5)] \textbf{$\F_1$-FST and the additive-multiplicative
  coupling.}  The (L1-K) obstruction of
  Section~\ref{sec:emergence} coincides with the classical
  Connes-Consani arithmetic-site programme; this is the long-term
  frame in which Route~C could conceivably be resurrected.
\item[(OFQ6)] \textbf{Physical realisation.}  Bost-Connes in
  condensed matter; experimental phase transition (speculative).
\item[(OFQ7)] \textbf{Rigidity.}  Adelic structure: Margulis-type
  rigidity or moduli?  (exploratory).
\end{description}

% =====================================================
\section{Status, Honest Assessment, and Outlook}
\label{sec:status}
% =====================================================

\subsection{What this paper establishes}

\begin{enumerate}[leftmargin=*]
\item Axiomatic FST framework for spectrum-zero identification,
  with explicit separation of baseline and completeness axioms
  (Sections~\ref{sec:axioms}, \ref{sec:fixed-point}).
\item Identification of (FST-7) with three structurally equivalent
  forms (Section~\ref{sec:fst7}); the three forms are, in practice,
  three distinct technical entry points, not interchangeable proof
  strategies.
\item Three technical obstacles (O1)--(O3) on Route~A, catalogued
  with candidate construction sketches but \emph{not} reduced to
  standard techniques (Section~\ref{sec:wick-mellin-obstacles}).
\item Selberg compatibility check: a non-abelian Tannakian extension
  is compatible with Selberg's RH proof
  (Section~\ref{sec:selberg}); it does not re-derive the proof.
\item Classification of zeta-type families
  (Table~\ref{tab:classification-v2}).
\item A \textbf{conditional reduction blueprint} for the two CCM
  missing steps on Route~B
  (Blueprint~\ref{bp:fst-transfer}), together with a
  four-milestone agenda for Route~B
  (Milestones~\ref{ms:1}--\ref{ms:4}).
\item \textbf{(New in v0.8)} A \textbf{Three-Lemma Endgame}
  (Section~\ref{sec:three-lemma-endgame}) that conditionally resolves
  (MS2) via mass-based spectral cluster analysis: three independent
  lemmas (Lemmas~\ref{lem:secular}--\ref{lem:coercive}) yield the
  mass concentration bound
  (Theorem~\ref{thm:mass-concentration}) as a corollary of
  a standard quasimode-to-projector inequality.
\item A diagnostic reformulation of Route~C identifying the
  additive-multiplicative coupling (L1-K) as its core obstruction
  (Section~\ref{sec:emergence}).
\end{enumerate}

\subsection{What this paper does not establish}

\begin{enumerate}[leftmargin=*]
\item No unconditional new theorem about RH or GRH.
\item No derivation of (FST-7) from the baseline FST axioms alone;
  only a conjectural \emph{relocation} of the arithmetic input under
  additional selection axioms.
\item No resolution of the Wick-Mellin obstacles (O1)--(O3); (O3)'s
  trace-class claim from v0.6 is explicitly retracted
  (Remark~\ref{rem:o3-retraction}).
\item No unconditional proof of (MS1) (simple-even for $QW_\lambda$);
  this remains open.
\item No unconditional proof of (MS2); the Three-Lemma Endgame
  (Section~\ref{sec:three-lemma-endgame}) resolves (MS2)
  \emph{conditionally} on Hypothesis~\ref{hyp:galerkin}
  (Galerkin faithfulness).
\item No proof of Hypothesis~\ref{hyp:galerkin} itself; it is
  numerically verified but not rigorously established in the
  $\lambda \to \infty$ limit.
\item No proof of a generalised Weyl asymptotic for $QW_\lambda$;
  replaced by Milestone~\ref{ms:4}.
\item No proof of the non-abelian fixed-point theorem
  (Conjecture~\ref{conj:nonabelian-fixed-point}).
\item No proof of Proposition~\ref{prop:thermo-derivation} for (B3);
  no proof of Conjecture~\ref{conj:l1k} (L1-K); no resolution of
  the additive-multiplicative coupling.
\end{enumerate}

\subsection{Honest assessment}

\textbf{FST in its present form is a comparative framework with one
conditional endgame result.}  It provides structural clarity, a
classification criterion, and an explicit three-route picture in
which the hard open content is localised:

\begin{itemize}[leftmargin=*]
\item On Route~A, the hard content sits in the Wick-Mellin obstacles
  (O2), (O3).
\item On Route~B, the hard content has been \emph{reduced} by the
  Three-Lemma Endgame.  The remaining open items are: (i)~(MS1)
  (simple-even, fully open); (ii)~Hypothesis~\ref{hyp:galerkin}
  (Galerkin faithfulness in the $\lambda \to \infty$ limit, numerically
  verified but not proved).  With Hypothesis~\ref{hyp:galerkin}
  granted, (MS2) is closed and Route~B reduces to (MS1) alone.
\item On Route~C, the hard content sits in Conjecture~\ref{conj:l1k}
  (L1-K), which is the classical additive-multiplicative coupling
  problem and is outside the reach of the present project.
\end{itemize}

The project does not constitute a proof of RH.  However, the status
has materially advanced: Route~B's approximation problem (MS2),
which was previously identified as one of two deep open items in the
CCM programme, is now conditionally resolved.  The condition
(Hypothesis~\ref{hyp:galerkin}) is a standard-looking Galerkin
convergence statement that is numerically verified to 12~digits of
precision and is expected to follow from existing finite-element
theory for rank-one perturbations of self-adjoint operators.

\subsection{Role of the framework}

\begin{itemize}
\item \emph{Uniform language} for Tate, Meyer, Bost-Connes, Selberg,
  CCM, and Connes-Consani.
\item \emph{Classification criterion}
  (Table~\ref{tab:classification-v2}).
\item \emph{Research agenda} (OFQ1--OFQ7; milestones
  \ref{ms:1}--\ref{ms:4} for Route~B).
\item \emph{Conditional endgame} for Route~B's (MS2) via the
  Three-Lemma Endgame (Section~\ref{sec:three-lemma-endgame}).
\item \emph{Diagnostic identification} that Route~C's gap is the
  additive-multiplicative coupling.
\item \emph{Compatibility sketch} with Selberg
  (Section~\ref{sec:selberg}).
\end{itemize}

\subsection{Next technical steps}

\begin{enumerate}[leftmargin=*]
\item \textbf{Route B, Milestone~\ref{ms:1}:} Domain/form theory for
  $\Psi_\lambda$ as a bounded intertwiner between suitable form
  domains.  This is the first rigorous step and is the prerequisite
  for any precise statement of Blueprint~\ref{bp:fst-transfer}.
\item \textbf{Route B, Milestone~\ref{ms:2}:} min-max low-eigenvalue
  comparison, not full Weyl asymptotic.
\item \textbf{Route B, Milestone~\ref{ms:3}:} finite-dimensional CCM
  truncation theorem inside $E_N$.
\item \textbf{Route B, Milestone~\ref{ms:4}:} only after 1--3,
  asymptotic regime.
\item \textbf{Route A:} operator-theoretic analysis of (O3) and
  precise articulation of the limiting KMS states required by (O2);
  the trace-class question
  (Conjecture~\ref{conj:ren-trace-class}) is open.
\item \textbf{Route C:} documentation only; no further attack
  without progress on additive-multiplicative coupling.
\item \textbf{Non-abelian (OFQ3):}  Langlands functoriality may be a
  better long-term frame than the present Tannakian FST formulation;
  cf.~\cite{TannakianAxioms}.
\item \textbf{Axiomatics:} independence/minimality proofs for the
  baseline axioms; separation of setup, dynamical, and
  arithmetic-completeness assumptions
  (Remark~\ref{rem:axioms-status}).
\end{enumerate}

Each of these is a multi-month research task.

\subsection{Relation to the Math-Master}

The present paper's Tannakian extension (Section~\ref{sec:selberg})
refines the Math-Master HP-BL classification
\cite{MathMaster} \S\S2.4, 9: even on the NO side (Riemann,
Dirichlet), an FST-spectral realisation may exist via the
idele-side SGE-YES route.  The three-component decomposition
$\mathrm{RH}(X) \iff \mathrm{GBC}(X) \wedge C^{(2)}(X) \wedge
\mathrm{Hurwitz}(X)$ of the Math-Master gets its $C^{(2)}$
component identified with the Meyer construction (as discussed in
\cite{C2LayerStatus}).  This identification is structural, not a
proof of the Math-Master decomposition's implications.

% =====================================================
\section*{Acknowledgements}
% =====================================================

This paper is a position draft prepared within the META\_RH\_TREE
project of FST-Mathematics.  Relative to v0.6, it is an honest
rewrite that incorporates the April 2026 external audit
(\texttt{\_audit/GPT\_REVIEW\_2026-04-17\_*.md}) and brings
paper-level claims into line with the honest note-level diagnoses.
It consolidates the Session-11 to Session-15 explorations documented
in the project folder
\texttt{fst\_spectrum\_duality/}, including:
\texttt{KONZEPT.md},
\texttt{Exploration.md},
\texttt{FST\_FIXPOINT\_THEOREM\_ATTEMPT.md},
\texttt{MINIMAL\_EXAMPLE\_RIEMANN.md},
\texttt{OPTION\_A\_PONTRJAGIN\_ABSTRACTION.md},
\texttt{OPTION\_B1\_INFORMATION\_THEORY.md},
\texttt{OPTION\_B3\_ARITHMETIC\_ERGODICITY.md},
\texttt{BRIDGE\_BOST\_CONNES\_MEYER.md},
\texttt{WICK\_ROTATION\_TECHNIQUE.md},
\texttt{O1\_DISTRIBUTION\_CLASS.md},
\texttt{O2\_FOURIER\_DUALITY.md},
\texttt{O3\_CONVERGENCE\_CONTROL.md},
\texttt{SELBERG\_CONSISTENCY\_CHECK.md},
\texttt{EMERGENCE\_HYPOTHESIS\_ATTACK.md},
\texttt{FREE\_ENERGY\_PRODUCT\_FORMULA.md},
\texttt{LEMMA\_3\_2\_PROOF.md},
\texttt{THERMODYNAMIC\_DERIVATION.md},
\texttt{TANNAKIAN\_FST\_AXIOMS.md},
\texttt{L1\_VARIATIONAL\_PROOF.md},
\texttt{L1\_K\_ATTACK.md},
\texttt{CCM\_2025\_STUDY.md},
\texttt{QW\_PW\_ISOMETRY.md}.

% =====================================================
\begin{thebibliography}{99}

\bibitem{MathMaster}
L.~Geiger,
\textit{The Zeta Zoo --- The Mathematical Side of Functional
Stability Theory},
Preprint (2026),
\texttt{00\_MASTERS/math\_master/paper/NE\_B\_BOUNDARY\_v1\_en.tex}.

\bibitem{RFEPFramework}
L.~Geiger,
\textit{Functional Stability Theory --- Physical Instantiation via
the Renormalized Free Energy Principle},
Preprint (2026) (v1.8 in preparation).

\bibitem{Geiger2026RHII}
L.~Geiger,
\textit{Riemann Hypothesis via Direct Frontier Dominance
(Paper II: Even Dominance)},
Preprint (2026),
Zenodo DOI \href{https://doi.org/10.5281/zenodo.19546593}{10.5281/zenodo.19546593}.

\bibitem{AnalyticKernelV2}
L.~Geiger,
\textit{Analytic Kernel V2 --- Rigorous Weil-Formula Derivation
of the Sign Coefficient},
Supplementary note (2026),
\texttt{02\_FST\_MATHEMATICS/dirichlet\_atlas/ANALYTIC\_KERNEL\_V2.md}.

\bibitem{C2LayerStatus}
L.~Geiger,
\textit{$C^{(2)}$-Layer Status Dirichlet},
Supplementary note (2026),
\texttt{02\_FST\_MATHEMATICS/dirichlet\_atlas/C2\_LAYER\_STATUS\_2026-04-17.md}.

\bibitem{Meyer2005}
R.~Meyer,
\textit{On a representation of the idele class group related to
primes and zeros of $L$-functions},
Duke Math. J. \textbf{127} (2005), 519--595.

\bibitem{Connes1999}
A.~Connes,
\textit{Trace formula in noncommutative geometry and the zeros of
the Riemann zeta function},
Selecta Math. (N.S.) \textbf{5} (1999), 29--106.

\bibitem{Connes2026}
A.~Connes,
\textit{Recent developments in the Riemann Hypothesis programme},
arXiv:2602.04022 (2026).

\bibitem{ConnesConsaniMoscovici2025}
A.~Connes, C.~Consani, H.~Moscovici,
\textit{Zeta Spectral Triples},
arXiv:2511.22755 (2025).

\bibitem{ConnesMoscovici2022}
A.~Connes, H.~Moscovici,
\textit{The UV prolate spectrum matches the zeros of zeta},
Proc. Natl. Acad. Sci. USA \textbf{119} (22) (2022), e2123174119;
arXiv:2112.05500.

\bibitem{ConnesConsani2023}
A.~Connes, C.~Consani,
\textit{Spectral triples and $\zeta$-cycles},
Enseign. Math. \textbf{69} (2023), 93--148; arXiv:2106.01715.

\bibitem{SlepianPollack1961}
D.~Slepian, H.~O.~Pollack,
\textit{Prolate spheroidal wave functions, Fourier analysis and
uncertainty --- I},
Bell Syst. Tech. J. \textbf{40} (1961), 43--63.

\bibitem{ConnesConsaniMoscovici2024}
A.~Connes, C.~Consani, H.~Moscovici,
\textit{Zeta zeros and prolate wave operators},
arXiv:2310.18423 (2024).

\bibitem{ConnesConsaniWeilPositivity2020}
A.~Connes, C.~Consani,
\textit{Weil positivity and trace formula --- the archimedean place},
Selecta Math. (N.S.) \textbf{27} (2021), art.~77; arXiv:2006.13771.

\bibitem{Kato}
T.~Kato,
\textit{Perturbation Theory for Linear Operators},
Grundlehren der math.\ Wiss.\ \textbf{132}, Springer (1966; 2nd ed.\ 1980).

\bibitem{L1KAttack}
L.~Geiger,
\textit{L1-K --- Attack on the Reconstruction Problem},
Supplementary note (2026),
\texttt{02\_FST\_MATHEMATICS/fst\_spectrum\_duality/L1\_K\_ATTACK.md}.

\bibitem{CCMStudy}
L.~Geiger,
\textit{CCM 2025 Zeta Spectral Triples --- Study + FST connection},
Supplementary note (2026),
\texttt{02\_FST\_MATHEMATICS/fst\_spectrum\_duality/CCM\_2025\_STUDY.md}.

\bibitem{BostConnes1995}
J.-B.~Bost, A.~Connes,
\textit{Hecke algebras, type~III factors and phase transitions
with spontaneous symmetry breaking in number theory},
Selecta Math. (N.S.) \textbf{1} (1995), 411--457.

\bibitem{Tate1950}
J.~T.~Tate,
\textit{Fourier analysis in number fields and Hecke's
zeta-functions},
Ph.D. thesis, Princeton University (1950). Reprinted in
Cassels-Fr\"ohlich, \emph{Algebraic Number Theory}, 1967.

\bibitem{CasselsFrohlich}
J.~W.~S.~Cassels, A.~Fr\"ohlich (eds.),
\textit{Algebraic Number Theory},
Academic Press (1967).

\bibitem{ConnesMarcolli}
A.~Connes, M.~Marcolli,
\textit{Noncommutative Geometry, Quantum Fields, and Motives},
American Mathematical Society Colloquium Publications 55 (2008).

\bibitem{Selberg1956}
A.~Selberg,
\textit{Harmonic analysis and discontinuous groups in weakly
symmetric Riemannian spaces with applications to Dirichlet series},
J. Indian Math. Soc. (N.S.) \textbf{20} (1956), 47--87.

\bibitem{CornelissenMarcolli2010}
G.~Cornelissen, M.~Marcolli,
\textit{Quantum Statistical Mechanics, L-series and Anabelian
Geometry I: Partition Functions},
arXiv:1009.0736 (2010).

\bibitem{Laca1998}
M.~Laca,
\textit{Semigroups of $\ast$-endomorphisms, Dirichlet series, and
phase transitions},
J. Funct. Anal. \textbf{152} (1998), 330--378.

\bibitem{LacaRaeburn2010}
M.~Laca, I.~Raeburn,
\textit{Phase transition on the Toeplitz algebra of the affine
semigroup over the natural numbers},
Adv. Math. \textbf{225} (2010), 643--688.

\bibitem{CuntzDeningerLaca2013}
J.~Cuntz, C.~Deninger, M.~Laca,
\textit{$C^\ast$-algebras of Toeplitz type associated with
algebraic number fields},
Math. Ann. \textbf{355} (2013), 1383--1423.

\bibitem{ArtinWhaples1945}
E.~Artin, G.~Whaples,
\textit{Axiomatic characterization of fields by the product formula
for valuations},
Bull. Amer. Math. Soc. \textbf{51} (1945), 469--492.

\bibitem{EmergenceAttack}
L.~Geiger,
\textit{Emergence Hypothesis --- Attack Document},
Supplementary note (2026),
\texttt{02\_FST\_MATHEMATICS/fst\_spectrum\_duality/EMERGENCE\_HYPOTHESIS\_ATTACK.md}.

\bibitem{LemmaThreeTwo}
L.~Geiger,
\textit{Lemma 3.2 --- $\N^\times$ as the unique FST$_0$-semigroup},
Supplementary note (2026),
\texttt{02\_FST\_MATHEMATICS/fst\_spectrum\_duality/LEMMA\_3\_2\_PROOF.md}.

\bibitem{ThermoDerivation}
L.~Geiger,
\textit{Thermodynamic Derivation of (B1)--(B4) and FST-Minimality},
Supplementary note (2026),
\texttt{02\_FST\_MATHEMATICS/fst\_spectrum\_duality/THERMODYNAMIC\_DERIVATION.md}.

\bibitem{TannakianAxioms}
L.~Geiger,
\textit{Tannakian FST Axioms --- Non-abelian Extension},
Supplementary note (2026),
\texttt{02\_FST\_MATHEMATICS/fst\_spectrum\_duality/TANNAKIAN\_FST\_AXIOMS.md}.

\end{thebibliography}

\end{document}
